% c00.fig.leiter — the ladder of theories of the course, from zoomed in to zoomed out.
%
% A scheme, not a plot.  Content: docs/plan/BRIEF.md:40 (TJO's wording, quoted in the
% English labels are selected excerpts; the full final qualifier is in the caption.
% The German labels are the chapter author's translation,
% gap G-0022).  Six rungs, bottom = zoomed in, top = zoomed out; five arrows, each
% labelled with the brief's words for that step ("zoom out", once "zoom out and lose
% info").  The rungs are numbered 1..6 so that the picture does not depend on colour.
\begin{tikzpicture}[x=1cm, y=1cm]
  \foreach \k/\y in {1/0, 2/1.25, 3/2.5, 4/3.75, 5/5, 6/6.25} {
    \draw[tpc axis] (0,\y) rectangle (11.2,\y+0.75);
    \node[tpc label] at (0.38,\y+0.375) {\k};
    \draw[tpc grid] (0.76,\y) -- (0.76,\y+0.75);
  }
  \node[tpc label, anchor=west] at (0.9,0.375)
    {\tpcifenglish{fundamental theory of physics (GPT = general probabilistic theory)}{Grundtheorie der Physik (GPT = allgemeine probabilistische Theorie)}};
  \node[tpc label, anchor=west] at (0.9,1.625)
    {\tpcifenglish{QM of single systems}{QM einzelner Systeme}};
  \node[tpc label, anchor=west] at (0.9,2.875)
    {\tpcifenglish{many particles}{viele Teilchen}};
  \node[tpc label, anchor=west] at (0.9,4.125)
    {\tpcifenglish{indefinite number of particles (Fock space)}{unbestimmte Teilchenzahl (Fockraum)}};
  \node[tpc label, anchor=west] at (0.9,5.375)
    {\tpcifenglish{stat mech: ensembles, fundamental hypothesis of stat mech}{statistische Mechanik: Ensembles, Grundhypothese der statistischen Mechanik}};
  \node[tpc label, anchor=west] at (0.9,6.625)
    {\tpcifenglish{thermodynamics}{Thermodynamik}};
  \foreach \y in {0.75, 2, 3.25, 5.75} {
    \draw[tpc axis, -{Stealth[length=5pt]}] (5.6,\y) -- (5.6,\y+0.5);
    \node[tpc label, anchor=west] at (5.75,\y+0.25) {\tpcifenglish{zoom out}{herauszoomen}};
  }
  \draw[tpc axis, -{Stealth[length=5pt]}] (5.6,4.5) -- (5.6,5);
  \node[tpc label, anchor=west] at (5.75,4.75)
    {\tpcifenglish{zoom out and lose info}{herauszoomen und Information verlieren}};
  \node[tpc label, anchor=east, rotate=90] at (-0.35,7.0) {\tpcifenglish{macroscopic}{makroskopisch}};
  \node[tpc label, anchor=west, rotate=90] at (-0.35,0) {\tpcifenglish{microscopic}{mikroskopisch}};
  \draw[tpc leader, -{Stealth[length=4pt]}] (-0.35,2.6) -- (-0.35,4.4);
\end{tikzpicture}
